\section{The Low-Period Spectral Frontier}

We prove Theorem F by combining complete orbit enumeration with a compressed
set of exact spectral exclusions. Its domain is the set of infinite operators
(2.1) whose Hamilton-gauge triangle word \(\tau\) has a least positive period
\(p\le 16\); the minimized quantity is the squared infinite-volume radius (2.5),
not a finite-holonomy spectrum. A word displayed in a repeated cell is kept
only for complete certificate accounting and is identified with its primitive
phase by zone folding. The theorem does not extrapolate beyond this domain.

\subsection{Complete phase space}

For each displayed cell length \(1\le p\le 16\), consider the legal flux words

\begin{lstlisting}[language={}]
Q in {+1,-1}^p,       product_i Q_i=1,                          (8.1)
\end{lstlisting}

modulo rotations and reflections. Explicit orbit partition and the independent
Burnside calculation of Appendix A give the counts

\begin{lstlisting}[language={}]
1,2,2,4,4,8,9,18,23,44,63,122,190,362,612,1162,                (8.2)
\end{lstlisting}

whose sum is 2,626. Every periodic phase of primitive \(\tau\) period at most 16
appears in the list when written in its primitive cell. Translation,
reflection, global edge-sign negation, and cell repetition have the spectral
equivalences proved in Lemmas 2.1-2.2.

\subsection{Exact exclusion scheme}

For every representative compute the moments (2.9). Lemma 2.3 supplies the
valid implication

\begin{lstlisting}[language={}]
F_k>0 ==> R(Q)>8>eta.                                            (8.3)
\end{lstlisting}

An adaptive exact closed-walk calculation through \(F_64\) finds a first
positive excess for 2,611 representatives. The first-positive indices range
from 1 to 64; in particular, two near-barrier classes are first detected only
at \(F_48\) and \(F_64\). This calculation uses integer closed-walk sums, so the
sign of every excess is exact.

Fifteen representatives remain. Eight are displayed even-period repetitions
of the all-negative phase. Equation (6.1) proves \(R=8>\eta\) for all of them.
Two are

\begin{lstlisting}[language={}]
p=8:  Q=00010001,
p=16: Q=0001000100010001,                                      (8.4)
\end{lstlisting}

where zero denotes negative flux and one positive flux. Lemma 2.2 shows that
the second row is the doubled-cell representation of the first, not a second
phase; Section 5 gives their common value \(R=\eta\).

The remaining five representatives are

\begin{center}
\begin{tabular}{ll}
\toprule
\(p\) & canonical \(Q\) \\
\midrule
10 & \(0000010001\) \\
12 & \(000000010001\) \\
14 & \(00000000010001\) \\
14 & \(00010010001001\) \\
16 & \(0000000000010001\) \\
\bottomrule
\end{tabular}%
\end{center}

For each row choose \(z \in {+1,-1}\) and a stored nonzero integer vector \(v\).
Let \(H=H_Q(z)\) and

\begin{lstlisting}[language={}]
r=(v^T H^2 v)/(v^T v).                                         (8.5)
\end{lstlisting}

All arithmetic in (8.5) is integral before the final rational division. The
five certificates first verify \(r>4\). Put

\begin{lstlisting}[language={}]
u=((r-4)^2-10)/2.
\end{lstlisting}

They then verify \(u>0\) and \(u^2>5\). Since \(r-4>0\), these inequalities select
the positive square-root branch and give

\begin{lstlisting}[language={}]
(r-4)^2>10+2sqrt(5),
r>4+sqrt(10+2sqrt(5))=eta.                                     (8.6)
\end{lstlisting}

The condition \(r>4\) is indispensable: the two inequalities involving \(u\)
alone can also hold on the lower branch. Rayleigh's principle and (8.5)-(8.6)
give \(R(Q)\ge r>\eta\) for each residual representative.

\subsection{Completeness accounting}

The exact partition is

\begin{lstlisting}[language={}]
2611  moment-detected representatives,
   8  repeated all-negative representatives,
   5  endpoint-certified residual competitors,
   2  displayed target representations,
----
2626  total representatives.                                   (8.7)
\end{lstlisting}

The sets in (8.7) are disjoint, and Appendix A proves that the 2,626
representatives are the full legal orbit space. Every non-target phase in the
domain therefore has \(R(Q)>\eta\), while the two target rows are one primitive
period-eight phase with \(R(Q)=\eta\). This proves Theorem F. \(\square\)

The nearest observed competitor occurs at period ten, but numerical ordering
within a non-target period is not part of the theorem. What is proved exactly
is the strict comparison of every competitor with \(\eta\).
